% CAPÍTULO 2 - FUNDAMENTAÇÃO TEÓRICA
%
% Edite o conteúdo com sua fundamentação. Use \section{} e \subsection{}
% para organizar os tópicos. Abaixo há exemplos de equação, figura e
% referência cruzada que você pode usar como modelo.

\chapter{FUNDAMENTAÇÃO TEÓRICA}
\label{cap:fundamentacao}

[Texto introdutório do capítulo.]

% -----------------------------------------------------------------------------
\section{[Primeiro grande tópico]}
\label{sec:topico-1}

\subsection{Definições e contexto}
\label{sec:topico-1-def}

[Apresente os conceitos fundamentais.]

\subsection{Propriedades relevantes}
\label{sec:topico-1-prop}

[Discuta as propriedades pertinentes ao trabalho.]

\subsection{Aplicações}
\label{sec:topico-1-apl}

[Descreva as aplicações relacionadas.]

% -----------------------------------------------------------------------------
\section{[Segundo grande tópico]}
\label{sec:topico-2}

[Exemplo de equação numerada com \texttt{\textbackslash label}:]

\begin{equation}
    \Delta G = \Delta E + \Delta \mathrm{ZPE} - T\,\Delta S,
    \label{eq:deltaG}
\end{equation}
onde $\Delta E$ é a variação de energia eletrônica, $\Delta\mathrm{ZPE}$
é a correção de energia de ponto zero, $T$ é a temperatura e $\Delta S$
é a variação de entropia.

[Exemplo de figura:\footnote{Descomente as linhas abaixo e coloque a
imagem em \texttt{figuras/}.}]

% \begin{figure}[htb]
%     \centering
%     \caption{Legenda da figura.}
%     \includegraphics[width=0.7\textwidth]{figuras/exemplo}
%     \label{fig:exemplo}
%     \fonte{Elaborado pelo autor.}
% \end{figure}

% -----------------------------------------------------------------------------
\section{[Terceiro grande tópico]}
\label{sec:topico-3}

[Continue conforme necessário.]

% -----------------------------------------------------------------------------
\section{Trabalhos correlatos e lacuna de pesquisa}
\label{sec:correlatos}

\subsection{[Subgrupo de trabalhos]}
\label{sec:correlatos-1}

[Descreva trabalhos relacionados.]

\subsection{[Outro subgrupo de trabalhos]}
\label{sec:correlatos-2}

[Descreva outros trabalhos.]

\subsection{Síntese e identificação da lacuna}
\label{sec:lacuna}

[Sintetize os trabalhos revisados e identifique a lacuna que
este trabalho pretende preencher.]
